% !TEX root = ../main.tex
\section{Introduction}

Fault localization (FL) is a fundamental step in automated debugging and program repair. Its goal is to identify the program locations that are most likely responsible for an observed failure, thereby reducing the search space for developers and downstream repair systems. Over the past decades, FL techniques have evolved from spectrum-based and mutation-based suspiciousness ranking to learning-based approaches that exploit code, tests, execution traces, and bug reports. More recently, large language models (LLMs) have introduced a new paradigm for FL: instead of only computing suspiciousness scores from predefined execution signals, an LLM-based localization agent can read natural-language reports, search a repository, inspect code, and reason about candidate functions. This shift makes FL not only a ranking problem, but also an evidence-seeking process over large and heterogeneous codebases.

Most existing FL studies are still grounded in relatively controlled and test-centered settings such as Defects4J, where bugs are reproducible and are accompanied by test suites and fault-triggering tests [REF]. Such settings naturally support FL techniques that rely on test execution signals, coverage information, and suspiciousness scores. However, modern repository-level LLM agents are increasingly evaluated in a different setting: given a natural-language issue and a full codebase, the agent must infer the relevant files and functions before any repair can succeed. SWE-bench exemplifies this issue-driven repository-level setting: each task is constructed from a real GitHub issue and its corresponding pull request, and the model is asked to edit the codebase to resolve the issue [REF]. Therefore, FL in this setting becomes more open-ended. The agent may need to understand project-specific abstractions, distinguish root-cause boundaries from symptom reporters, compare working and failing code paths, and avoid misleading keyword matches.

Despite this richer interaction process, most FL methods still treat each bug as an isolated localization instance. Once a suspiciousness ranking or a list of candidate functions is produced, the intermediate localization trajectory is usually discarded. This trajectory may contain valuable reusable experience: which repository boundary turned out to be decisive, which searched locations were misleading, which candidate functions merely reported the symptom, and which reasoning pattern helped distinguish the root cause. In SWE-bench-style FL, such experience is especially important because the agent must repeatedly solve repository-level localization problems without the dense coverage signals assumed by many traditional FL settings. This motivates a new question: can historical localization trajectories be transformed into an explicit memory that improves future FL cases?

Recent work on agent skills offers a promising direction. A skill can serve as an external natural-language state that adapts a frozen agent without modifying its weights. For example, SkillOpt treats the skill document as a trainable object and uses an optimizer model to convert scored rollouts into bounded add/delete/replace edits; its reflection process uses failure trajectories to propose missing or corrective rules and success trajectories to preserve behaviors that already work [REF]. This line of work suggests that agent experience can be accumulated as editable textual memory rather than being discarded after each run.

However, fault localization introduces a domain-specific attribution problem that is not captured by a single final success or failure signal. An FL agent does not merely produce an answer; it interprets the issue, searches the repository, inspects candidate files and functions, relates symptoms to code behavior, distinguishes symptom reporters from root-cause boundaries, and finally ranks suspicious functions. A Top-5 miss may therefore result from different kinds of missing experience: the agent may misunderstand the project structure, apply the wrong fault mechanism, follow misleading evidence, or rank a downstream observer above the true causal boundary. Conversely, a Top-5 hit may depend on only one particular kind of useful experience, while other retrieved knowledge is incidental. Thus, directly optimizing a monolithic skill document from the final hit/miss outcome risks mixing heterogeneous localization lessons and makes it unclear which part of the skill memory should be edited.

To address this problem, we propose a layered view of self-evolving FL skills. Instead of treating a skill as a single procedural workflow, we represent historical localization experience as dimension-level knowledge skills. Different dimensions capture different kinds of reusable FL knowledge, such as project structure, fault mechanisms, evidence-interpretation strategies, and general localization principles. At runtime, relevant skills are retrieved and assembled into a source-marked localization context for a frozen Explorer. During evolution, both successful and failed trajectories can be used as feedback, but each bounded edit must target a specific source skill block or a missing-dimension slot. In this way, useful localization experience can be preserved, misleading knowledge can be refined, and new skills can be created only when a trajectory reveals transferable knowledge that is absent from the current SkillBank.
